\rSec3[optional.ref.monadic]{Monadic operations}

\indexlibrarymember{and_then}{optional}%
\begin{itemdecl}
template<class F>
  constexpr auto and_then(F&& f) const;
\end{itemdecl}

\begin{itemdescr}
\pnum
Let \tcode{U} be \tcode{invoke_result_t<F, T\&>}.

\pnum
\mandates
\tcode{U} is a specialization of \tcode{optional}.

\pnum
\effects
Equivalent to:
\begin{codeblock}
using U = invoke_result_t<F, T&>;

if (has_value()) {
  return invoke(forward<F>(f), *@\exposid{val}@);
} else {
  return remove_cvref_t<U>();
}
\end{codeblock}
\end{itemdescr}

\indexlibrarymember{transform}{optional}%
\begin{itemdecl}
template<class F>
  constexpr optional<invoke_result_t<F, T&>> transform(F&& f) const;
\end{itemdecl}

\begin{itemdescr}
\pnum
Let \tcode{U} be \tcode{invoke_result_t<F, T\&>}.

\pnum
\mandates
\tcode{U} is a valid contained type for \tcode{optional}. The declaration
\begin{codeblock}
U u(invoke(std::forward<F>(f), *@\exposid{val}@));
\end{codeblock}
is well-formed for some invented variable \tcode{u}.
\begin{note}
There is no requirement that \tcode{U} is movable\iref{dcl.init.general}.
\end{note}

\pnum
\returns
If \tcode{*this} contains a value, an \tcode{optional<U>} object whose contained value is direct-non-list-initialized with \tcode{invoke(std::forward<F>(f), *\exposid{val})}; otherwise, \tcode{optional<U>()}.
\end{itemdescr}

\indexlibrarymember{or_else}{optional}%
\begin{itemdecl}
template<class F>
  constexpr optional or_else(F&& f) const;
\end{itemdecl}

\begin{itemdescr}
\pnum
\constraints
\tcode{F} models \libconcept{invocable}.

\pnum
\mandates
\tcode{is_same_v<remove_cvref_t<invoke_result_t<F>>, optional>} is \tcode{true}.

\pnum
\effects
Equivalent to:
\begin{codeblock}
if (has_value()) {
  return *this;
} else {
  return forward<F>(f)();
}
\end{codeblock}
\end{itemdescr}
